% ============================================================================
%  styles/cover.tex — 出版家具：封面 / 扉页 / 版权页 / 封底 / 卷隔页 / 章扉页
%
%  全部版式随 \paperwidth / \paperheight 相对定位，换开本无需改数值；
%  颜色全部来自主题 token，换肤即时生效。
%
%  实现机制（书籍化改造的关键）：
%    全出血页 = 普通页（无 \newgeometry）+ 正文流内的
%    tikzpicture[remember picture, overlay]（\LT@FullPage 壳）。
%    不要把带透明度的绘图搬进 shipout 钩子（eso-pic / l3 shipout）：
%    本工具链下钩子材料里的 ExtGState 资源登记不可靠，pgf 的 opacity
%    会被按实色渲染（淡色水印/淡洗会变成实色墨块）；正文流内则一切正常。
%    remember picture 坐标按惯例第二遍收敛（目录/交叉引用本来就要多遍）。
%
%  与原笔记模板的差别（书籍化改造的核心）：
%    · 删除「课程名 / 学生姓名 / 学号 / 日期」信息表 —— 那是作业报告家具，
%      出版形制换成「丛书行 / 卷次行 / 作者 / 出版方 / 版次 / 书号」。
%    · 删除参考文献栏与致谢栏 —— 移到前装致谢页。
%    · 封面增加左缘三色出血带与底部三色带，作为孤星母题的书籍化收敛。
% ============================================================================

% --- 全出血页壳：翻页 → 流内整页 overlay → 空盒撑页 → 付印 -------------------
% 坐标系：current page.*（remember picture）；\ifLT@fullbleedpage 供
% background 钩子跳过网格与三色角标（出血页自带母题）。
\newcommand{\LT@FullPage}[1]{%
  \clearpage
  \LT@fullbleedpagetrue
  \thispagestyle{empty}%
  \begin{tikzpicture}[remember picture, overlay]#1\end{tikzpicture}%
  \null
  \clearpage
  \LT@fullbleedpagefalse
}

% --- 左缘三色竖带（书脊侧，出血进入裁切区）----------------------------------
\newcommand{\LT@SpineBand}{%
  \fill[lv] (current page.south west) rectangle ++(0.011\paperwidth,\paperheight);
  \fill[huang] ([xshift=0.011\paperwidth]current page.south west)
    rectangle ++(0.011\paperwidth,\paperheight);
  \fill[hong] ([xshift=0.022\paperwidth]current page.south west)
    rectangle ++(0.011\paperwidth,\paperheight);%
}

% --- 内框细线（封面/隔页/章扉共用）------------------------------------------
\newcommand{\LT@HairFrame}{%
  \draw[InkMuted, opacity=0.4, line width=0.6pt]
    ([xshift=0.075\paperwidth, yshift=0.030\paperheight]current page.south west)
    rectangle
    ([xshift=-0.055\paperwidth, yshift=-0.030\paperheight]current page.north east);%
}

% --- 底部三色横带 ------------------------------------------------------------
\newcommand{\LT@TriBand}{%
  \fill[lv]   (current page.south west) rectangle ++(0.34\paperwidth,0.022\paperheight);
  \fill[huang] ([xshift=0.34\paperwidth]current page.south west)
    rectangle ++(0.33\paperwidth,0.022\paperheight);
  \fill[hong] ([xshift=0.67\paperwidth]current page.south west)
    rectangle ++(0.33\paperwidth,0.022\paperheight);%
}

% ============================================================================
%  封面
% ============================================================================
\newcommand{\FrontCover}{%
  \LT@FullPage{%
    % --- 纸面 ---
    \fill[PaperIvory] (current page.south west) rectangle (current page.north east);
    \LT@SpineBand

    % --- 极淡数字母题水印 ---
    \node[anchor=center, opacity=0.055, text=InkBrown]
      at ([xshift=0.06\paperwidth, yshift=0.16\paperheight]current page.center)
      {\resizebox{!}{0.30\paperheight}{\LT@watermark}};

    % --- 内框细线 ---
    \LT@HairFrame

    % --- 丛书行 ---
    \node[anchor=north west, text=InkMuted, inner sep=0pt]
      at ([xshift=0.10\paperwidth, yshift=-0.075\paperheight]current page.north west)
      {\kaishu\fontsize{12}{15}\selectfont\LT@series};
    \draw[InkBrown, line width=0.9pt]
      ([xshift=0.10\paperwidth, yshift=-0.100\paperheight]current page.north west) --
      ([xshift=0.62\paperwidth, yshift=-0.100\paperheight]current page.north west);

    % --- 主标题 ---
    \node[anchor=north west, text=InkBrown, inner sep=0pt,
          text width=0.72\paperwidth, align=left]
      at ([xshift=0.10\paperwidth, yshift=-0.32\paperheight]current page.north west)
      {\songti\bfseries\fontsize{34}{42}\selectfont\LT@title};
    \node[anchor=north west, text=InkMuted, inner sep=0pt,
          text width=0.72\paperwidth, align=left]
      at ([xshift=0.10\paperwidth, yshift=-0.455\paperheight]current page.north west)
      {\itshape\fontsize{12}{16}\selectfont\LT@subtitle};

    % --- 卷次行 ---
    \node[anchor=north west, text=structurecolor, inner sep=0pt]
      at ([xshift=0.10\paperwidth, yshift=-0.585\paperheight]current page.north west)
      {\kaishu\fontsize{15}{19}\selectfont\LT@volumeline};
    \node[anchor=north west, text=InkMuted, inner sep=0pt]
      at ([xshift=0.10\paperwidth, yshift=-0.635\paperheight]current page.north west)
      {\LoneDisplay\fontsize{9.5}{12}\selectfont\LT@volumelineEN};

    % --- 代表公式槽（图灵 LADR 风格）---
    \node[anchor=north west, text=InkMuted, inner sep=0pt]
      at ([xshift=0.10\paperwidth, yshift=-0.700\paperheight]current page.north west)
      {\fontsize{13}{18}\selectfont\LT@coverline};

    % --- 底部：作者 + 出版方 ---
    \draw[InkMuted, line width=0.6pt, opacity=0.75]
      ([xshift=0.10\paperwidth, yshift=0.105\paperheight]current page.south west) --
      ([xshift=0.90\paperwidth, yshift=0.105\paperheight]current page.south west);
    \node[anchor=south west, text=InkBrown, inner sep=0pt]
      at ([xshift=0.10\paperwidth, yshift=0.062\paperheight]current page.south west)
      {\kaishu\fontsize{14}{18}\selectfont\LT@author};
    \node[anchor=south east, text=InkMuted, inner sep=0pt]
      at ([xshift=-0.06\paperwidth, yshift=0.062\paperheight]current page.south east)
      {\kaishu\fontsize{11}{14}\selectfont\LT@publisher};

    % --- 底部三色横带 ---
    \LT@TriBand
  }%
}

% ============================================================================
%  扉页（环衬后）：丛书名 + 书名，安静地落在右页上
% ============================================================================
\newcommand{\HalfTitle}{%
  \LT@FullPage{%
    \fill[PaperIvory] (current page.south west) rectangle (current page.north east);
    \node[anchor=north west, text=InkMuted, inner sep=0pt]
      at ([xshift=0.13\paperwidth, yshift=-0.16\paperheight]current page.north west)
      {\kaishu\fontsize{12}{16}\selectfont\LT@series};
    \node[anchor=north west, text=InkBrown, inner sep=0pt,
          text width=0.70\paperwidth, align=left]
      at ([xshift=0.13\paperwidth, yshift=-0.44\paperheight]current page.north west)
      {\songti\bfseries\fontsize{30}{38}\selectfont\LT@title};
    \node[anchor=north west, text=InkMuted, inner sep=0pt]
      at ([xshift=0.13\paperwidth, yshift=-0.56\paperheight]current page.north west)
      {\itshape\fontsize{12}{16}\selectfont\LT@subtitle};
    \draw[structurecolor, line width=1.1pt]
      ([xshift=0.13\paperwidth, yshift=-0.60\paperheight]current page.north west) --
      ++(0.10\paperwidth,0);
    \node[anchor=south west, text=InkMuted, inner sep=0pt]
      at ([xshift=0.13\paperwidth, yshift=0.07\paperheight]current page.south west)
      {\kaishu\fontsize{10}{13}\selectfont\LT@publisher};
  }%
}

% ============================================================================
%  版权页：出版形制必备，取代原模板的「学生 / 学号」块
% ============================================================================
\newcommand{\Imprint}{%
  \LT@FullPage{%
    \fill[PaperIvory] (current page.south west) rectangle (current page.north east);
    \node[anchor=north west, text=InkMuted, inner sep=0pt,
          text width=0.62\paperwidth, align=left]
      at ([xshift=0.13\paperwidth, yshift=-0.14\paperheight]current page.north west)
      {%
        {\kaishu\fontsize{11}{19}\selectfont
          \textbf{\LT@title}\par
          \ifdefempty{\LT@subtitle}{}{\LT@subtitle\par}
          \vspace{1.2em}%
          \LT@author\par
          \vspace{2.2em}%
          \ifdefempty{\LT@isbn}{}{书号：\LT@isbn\par}
          \ifdefempty{\LT@edition}{}{版次：\LT@edition\par}
          出版：\LT@publisher\par
          \vspace{2.2em}%
          \textcopyright\ \the\year\ \LT@author\par
          保留所有权利。未经书面许可，不得以任何形式复制、转载或倒卖本书。\par
          \vspace{1.2em}%
          本书为书稿版本，内容与版式可能在作者后续修订中调整。\par
        }%
      };
    \node[anchor=south east, text=InkMuted, inner sep=0pt]
      at ([xshift=-0.075\paperwidth, yshift=0.06\paperheight]current page.south east)
      {\LoneDisplay\fontsize{9}{11}\selectfont Lone Trail Book Template};
  }%
}

% ============================================================================
%  封底：内容简介 + 分卷列表 + 出版标 + 条码位
% ============================================================================
\newcommand{\BackCover}{%
  \LT@FullPage{%
    \fill[PaperIvory] (current page.south west) rectangle (current page.north east);
    \node[anchor=center, opacity=0.05, text=InkBrown]
      at ([xshift=0.05\paperwidth, yshift=0.2\paperheight]current page.center)
      {\resizebox{!}{0.34\paperheight}{\LT@watermark}};

    \node[anchor=north west, text=InkBrown, inner sep=0pt,
          text width=0.62\paperwidth, align=left]
      at ([xshift=0.10\paperwidth, yshift=-0.12\paperheight]current page.north west)
      {%
        {\kaishu\bfseries\fontsize{16}{21}\selectfont 内容简介}\\[5mm]
        {\kaishu\fontsize{11.5}{19}\selectfont\color{InkMuted}%
          \LT@blurb
          \ifdefempty{\LT@shelflist}{}{\par\vspace{6mm}\LT@shelflist}}%
      };

    % 出版标 + 条码占位
    \node[anchor=south west, text=InkBrown, inner sep=0pt]
      at ([xshift=0.10\paperwidth, yshift=0.075\paperheight]current page.south west)
      {\kaishu\fontsize{13}{17}\selectfont\LT@publisher};
    \node[anchor=south east, text=InkMuted, inner sep=0pt]
      at ([xshift=-0.06\paperwidth, yshift=0.075\paperheight]current page.south east)
      {\kaishu\fontsize{10}{13}\selectfont\ifdefempty{\LT@isbn}{}{\LT@isbn}};

    \LT@TriBand
  }%
}

% ============================================================================
%  卷隔页（多卷合订本）
%  槽位宏：书籍工程覆写下列三个，即可换上自己的卷次数据（不改本文件）。
%    \LoneVolumeTitle{N}  卷次大标题（如「第一卷 · 基础语言」）
%    \LoneVolumeBlurb{N}  卷次导语
%    \LoneVolumeMark{N}   右上角卷标（默认「第N卷」）
% ============================================================================
\newcommand{\LoneVolumeTitle}[1]{}
\newcommand{\LoneVolumeBlurb}[1]{}
\newcommand{\LoneVolumeMark}[1]{第#1卷}

\newcommand{\VolumeDivider}[1]{% #1 = 卷号
  \LT@FullPage{%
    \fill[PaperIvory] (current page.south west) rectangle (current page.north east);
    \node[anchor=center, opacity=0.05, text=InkBrown]
      at ([xshift=0.08\paperwidth, yshift=0.14\paperheight]current page.center)
      {\resizebox{!}{0.26\paperheight}{\LT@watermark}};
    \LT@HairFrame
    % 顶部：丛书行（左）+ 卷标（右），文字完整坐在横线之上
    \draw[InkBrown, line width=0.9pt]
      ([xshift=0.10\paperwidth, yshift=-0.085\paperheight]current page.north west) --
      ([xshift=0.90\paperwidth, yshift=-0.085\paperheight]current page.north east);
    \node[anchor=north west, text=InkMuted, inner sep=0pt]
      at ([xshift=0.10\paperwidth, yshift=-0.060\paperheight]current page.north west)
      {\kaishu\fontsize{11}{14}\selectfont\LT@series};
    \node[anchor=north east, text=InkMuted, inner sep=0pt]
      at ([xshift=-0.075\paperwidth, yshift=-0.060\paperheight]current page.north east)
      {\kaishu\fontsize{9.5}{12}\selectfont\LoneVolumeMark{#1}};
    % 卷次大标题 + 导语
    \node[anchor=north west, text=InkBrown, inner sep=0pt,
          text width=0.66\paperwidth, align=left]
      at ([xshift=0.10\paperwidth, yshift=-0.34\paperheight]current page.north west)
      {\songti\bfseries\fontsize{30}{38}\selectfont\LoneVolumeTitle{#1}};
    \draw[InkBrown, line width=1.15pt]
      ([xshift=0.10\paperwidth, yshift=-0.45\paperheight]current page.north west) --
      ++(0.14\paperwidth,0);
    \node[anchor=north west, text=InkBrown, inner sep=0pt,
          text width=0.66\paperwidth, align=left]
      at ([xshift=0.10\paperwidth, yshift=-0.50\paperheight]current page.north west)
      {\kaishu\fontsize{12}{20}\selectfont\LoneVolumeBlurb{#1}};
  }%
}

% ============================================================================
%  章扉页：正文每章前的整页
%  书籍工程可覆写下列槽位宏，不改本文件即得到自己的章节数据。
% ============================================================================
\newcommand{\LonePlateSymbol}[1]{$\cdot$}
\newcommand{\LonePlatePartCN}[1]{}
\newcommand{\LonePlateQuote}[1]{}
\newcommand{\LonePlateQuoteBy}[1]{}
\newcommand{\LonePlateImage}[1]{}   % 返回图片路径，空则不放

\newcommand{\ChapterPlate}[2]{% #1 = 章号 #2 = 章名
  \LT@FullPage{%
    \fill[PaperIvory] (current page.south west) rectangle (current page.north east);

    % 右半幅可放数学家肖像（书籍工程提供）
    \edef\LT@plateimg{\LonePlateImage{#1}}%
    \ifdefempty{\LT@plateimg}{}{%
      \IfFileExists{\LT@plateimg}{%
        \begin{scope}[opacity=0.15]
          \node[anchor=south east, inner sep=0pt]
            at ([xshift=-0.02\paperwidth, yshift=0.07\paperheight]current page.south east)
            {\includegraphics[height=0.60\paperheight,keepaspectratio]{\LT@plateimg}};
        \end{scope}%
        % 左淡洗，保证标题可读
        \fill[PaperIvory, opacity=0.55]
          (current page.south west) rectangle
          ([xshift=0.50\paperwidth]current page.north west);
        \fill[PaperIvory, opacity=0.30]
          ([xshift=0.50\paperwidth]current page.south west) rectangle
          ([xshift=0.64\paperwidth]current page.north west);
      }{}%
    }%

    % 数学母题水印
    \node[anchor=center, opacity=0.06, text=InkMuted]
      at ([xshift=0.07\paperwidth, yshift=0.14\paperheight]current page.center)
      {\resizebox{!}{0.30\paperheight}{\LonePlateSymbol{#1}}};

    \LT@HairFrame

    % 顶部：丛书行 + 篇名（文字完整坐在横线之上）
    \draw[InkBrown, line width=0.9pt]
      ([xshift=0.10\paperwidth, yshift=-0.085\paperheight]current page.north west) --
      ([xshift=0.90\paperwidth, yshift=-0.085\paperheight]current page.north east);
    \node[anchor=north west, text=InkMuted, inner sep=0pt]
      at ([xshift=0.10\paperwidth, yshift=-0.060\paperheight]current page.north west)
      {\kaishu\fontsize{11}{14}\selectfont\LT@series};
    \node[anchor=north east, text=InkMuted, inner sep=0pt]
      at ([xshift=-0.075\paperwidth, yshift=-0.060\paperheight]current page.north east)
      {\kaishu\fontsize{9.5}{12}\selectfont\LonePlatePartCN{#1}};

    % 章号 + 章名
    \node[anchor=north west, text=InkMuted, inner sep=0pt]
      at ([xshift=0.10\paperwidth, yshift=-0.32\paperheight]current page.north west)
      {\kaishu\fontsize{14}{18}\selectfont 第\,#1\,章};
    \node[anchor=north west, text=InkBrown, inner sep=0pt,
          text width=0.66\paperwidth, align=left]
      at ([xshift=0.10\paperwidth, yshift=-0.37\paperheight]current page.north west)
      {\songti\bfseries\fontsize{30}{38}\selectfont #2};
    \draw[InkBrown, line width=1.15pt]
      ([xshift=0.10\paperwidth, yshift=-0.48\paperheight]current page.north west) --
      ++(0.14\paperwidth,0);

    % 题记
    \ifdefempty{\LonePlateQuote{#1}}{}{%
      \node[anchor=north west, text=InkBrown, inner sep=0pt,
            text width=0.62\paperwidth, align=left]
        at ([xshift=0.10\paperwidth, yshift=-0.54\paperheight]current page.north west)
        {%
          \begin{minipage}{0.62\paperwidth}
            {\kaishu\fontsize{12.5}{20}\selectfont
              「\,\LonePlateQuote{#1}\,」}\\[0.4cm]
            {\kaishu\fontsize{10}{13}\selectfont\color{InkMuted}
              --- \LonePlateQuoteBy{#1}}
          \end{minipage}%
        };
    }%

    \draw[InkMuted, line width=0.55pt, opacity=0.65]
      ([xshift=0.10\paperwidth, yshift=0.055\paperheight]current page.south west) --
      ([xshift=0.90\paperwidth, yshift=0.055\paperheight]current page.south west);
  }%
}
